\section{Model and notation}\label{sec:model}

\subsection{Bitcoin, sats and ordinal ranges}

A Bitcoin transaction consumes previous outputs and creates new ones. Each
output has a value in satoshis. Ordinal theory assigns the sats of the inputs
to the outputs in order, so an output can be described by a list of half-open
intervals $[a,b)$ of ordinal numbers. A coin's \emph{first sat} is the first
sat in its first interval. Transfers preserve this order: sats are spent
first-in-first-out, and a later output cannot jump ahead of an earlier one.

Two properties follow and are used throughout the protocol.

\begin{enumerate}[leftmargin=1.6em,itemsep=0.2em]
  \item A frog's anchor sat is the first sat of the output that carries it.
        In a mint, that is the first sat of the commit output; after a spend,
        it must be at offset zero of output zero.
  \item A coin that the indexer has never seen move has no true ordinal
        assignment. It receives a deterministic \emph{shadow} range so that
        FIFO placement remains well-defined, but a shadow range is not
        ownership. A mint whose anchor sat is a shadow is rejected as
        \code{MintAnchorUntracked}.
\end{enumerate}

The protocol never relies on a floating-point sat number or on a database
identifier. The sat is the identity, and the satpoint (transaction id and
output index) is its current location.

\subsection{Taproot and tapscript}

A Taproot output commits to an internal key $P$ and a Merkle root of a script
tree. The output key is $Q = P + tG$ for a tweak $t$ derived from the tree.
A key-path spend signs with the tweaked key; a script-path spend reveals one
leaf, the leaf script, and a control block that proves the leaf's position in
the tree. BIP-341 \cite{bip341} defines the output key and control block; BIP-342
\cite{bip342} defines the tapscript signature hash used by
\code{OP_CHECKSIG}.

HOPPER uses the script path exclusively. The internal key is the fixed
nothing-up-my-sleeve point \code{NUMS_XONLY}, for which no discrete log is
known. There is no key-path spend, ever. The tree has four leaves at depth
two: three \emph{cross} leaves, one per direction, each checking the owner's
signature and carrying the state envelope for that direction, and one
\emph{send} leaf that checks only the owner's signature. The exact scripts
and their branch structure are given in Section~\ref{sec:custody}.

\subsection{The ord envelope}

An ord inscription is a tapleaf containing an envelope of the form
\begin{center}
\code{OP_FALSE OP_IF "ord" (tag, value)... 0x00 (body chunks) OP_ENDIF}.
\end{center}
The protocol identifier is the three-byte push \code{"ord"}. Tag \code{1} is
the content type, tag \code{3} is a parent, and the empty push \code{0x00}
marks the start of the body. Body chunks are limited to 520 bytes and are
concatenated in order. The indexer's parser mirrors ord's decoder, including
ord's minimal-index rule for parent tags.

HOPPER writes only the fields it needs: a content type and a body, plus one
parent tag on a crossing. It never writes a pointer tag. A pointer would move
the inscription off the anchor sat, which would break the identity between
the frog and its state; a mint inscription carrying a pointer is rejected.

\subsection{Trust model}

The protocol distinguishes four roles.

\begin{table}[htbp]
\centering
\small
\caption{What each component is trusted to do.}
\label{tab:trust}
\begin{tabularx}{\linewidth}{@{}p{0.19\linewidth}X X@{}}
\toprule
\textbf{Component} & \textbf{Trusted for} & \textbf{Not trusted for} \\
\midrule
Bitcoin consensus & Transaction validity, ordering, block hashes, finality &
Application-level meaning \\
ord viewer & Parsing ord envelopes and serving \code{/r/} content &
Knowing HOPPER's rules \\
HOPPER indexer & Convenience: fast discovery, API, event log &
Custody, signatures, art truth, the outcome of a future block \\
Wallet & Key custody, transaction construction, user intent &
Chain history it has not checked \\
\bottomrule
\end{tabularx}
\end{table}

The art document is deliberately independent of all four. It embeds the fold
source and reads only relative endpoints on the viewer's own ord node. The
static band states only what the mint transaction can prove: the name, a
generation hint, and an unknown crossing count. It never prints a number.
The live fold may render a count only after it has verified every body-valid
child's raw transaction; a missing or false proof yields the
\emph{unverified} band, which prints no count and never claims \code{HOME}.
Only a settled frog gets \code{RIP}. If the viewer refuses every fetch or
runs with JavaScript disabled, the static band stays and does not assert that
the frog is alive. The indexer can be unavailable without changing what a
frog is.

The adversary is assumed to control network delivery, to be able to censor or
reorder transactions before confirmation, and to attempt to spend outputs it
does not own. It is not assumed to break secp256k1, SHA-256, BIP-340
signatures or Bitcoin consensus. Reorganizations are treated as a normal
chain event, not as an attack.

\subsection{Notation}

\begin{table}[htbp]
\centering
\small
\caption{Notation used in the remainder of the paper.}
\label{tab:notation}
\begin{tabularx}{\linewidth}{@{}p{0.20\linewidth}X@{}}
\toprule
\textbf{Symbol} & \textbf{Meaning} \\
\midrule
$G$ & Generation genesis height; the only deployment value \\
$\mathrm{sat}_g$ & The one grandfathered v7 sat, or $0$ when none \\
$b$ & Required crossings on the first block of a generation \\
$\mathit{gi}$ & Blocks per extra required crossing \\
$c$ & Ceiling, the hardest the road ever gets \\
$h_{\rm hold}$ & Blocks spent at the ceiling before the reset \\
$L$ & Generation length, $(c-b)\mathit{gi}+h_{\rm hold}$ \\
$S(k)$ & First height of generation $k$ (1-based) \\
$\gamma(h)$ & Generation number at height $h$ \\
$r(h)$ & Required crossings at height $h$ \\
$w(B)$ & The direction block $B$ wraps to, in $\{0,1,2\}$ \\
$T$ & Confirmation height of a crossing \\
$B$ & Outcome block, $B=T+1$ \\
$H_B$ & 32-byte hash of block $B$, internal byte order \\
$\kappa$ & Crossing finality, six blocks \\
$v$ & Wire version: \code{0x08}, or \code{0x07} for the grandfather \\
$\mathrm{sat}$ & The frog's ordinal number (a \code{u64}) \\
$\mathit{dir}$ & Chosen direction: $0$ left, $1$ straight, $2$ right \\
\bottomrule
\end{tabularx}
\end{table}

\noindent
All integer encodings in this paper are little-endian unless a field is
explicitly described as display-order hex. Transaction ids and block hashes
are shown to humans in reversed (display) order, as explorers do; the wrap
uses the internal order.
